\documentclass[11pt,a4paper]{article}
\usepackage[margin=24mm,headheight=15pt]{geometry}
\usepackage[T1]{fontenc}
\usepackage[utf8]{inputenc}
\usepackage{amsmath,amsthm}
\usepackage{newpxtext,newpxmath}
\usepackage{microtype,mathtools,bm}
\usepackage{booktabs,array,longtable,enumitem}
\usepackage[dvipsnames]{xcolor}
\usepackage{fancyhdr}
\usepackage[most]{tcolorbox}
\usepackage{hyperref}
\usepackage[nameinlink,noabbrev]{cleveref}
\hypersetup{colorlinks=true,linkcolor=MidnightBlue,citecolor=MidnightBlue,urlcolor=MidnightBlue,
 pdftitle={Certified Inversion Through the q-to-1 Transition},
 pdfauthor={Research manuscript prepared for Vladimir Reshetnikov}}
\newcommand{\ee}{\mathrm e}
\newcommand{\dd}{\,\mathrm d}
\newcommand{\Li}{\operatorname{Li}}
\newcommand{\E}{\mathcal E}
\newcommand{\D}{\mathcal D}
\newcommand{\R}{\mathbb R}
\newcommand{\N}{\mathbb N}
\newcommand{\Z}{\mathbb Z}
\newcommand{\supp}{\operatorname{supp}}
\newcommand{\sgn}{\operatorname{sgn}}
\newcommand{\Res}{\operatorname{Res}}
\newcommand{\bigO}{\mathcal O}
\newcommand{\qpoch}[2]{(#1;#2)_\infty}
\newcommand{\repo}{\textsc{ProveIt}}
\newcommand{\commit}{20301601562cf84173f2b503b0dc9c3fbf21cd0c}
\DeclareMathOperator{\diag}{diag}
\numberwithin{equation}{section}
\theoremstyle{plain}
\newtheorem{theorem}{Theorem}[section]
\newtheorem{proposition}[theorem]{Proposition}
\newtheorem{lemma}[theorem]{Lemma}
\newtheorem{corollary}[theorem]{Corollary}
\theoremstyle{definition}
\newtheorem{definition}[theorem]{Definition}
\newtheorem{example}[theorem]{Example}
\theoremstyle{remark}
\newtheorem{remark}[theorem]{Remark}
\newtcolorbox{resultbox}{colback=MidnightBlue!3,colframe=MidnightBlue!45,
 boxrule=.6pt,arc=1mm,left=3mm,right=3mm,top=2mm,bottom=2mm,breakable}
\pagestyle{fancy}
\fancyhf{}
\fancyhead[L]{\small Certified inversion through the $q\to1$ transition}
\fancyhead[R]{\small Research manuscript}
\fancyfoot[C]{\thepage}
\setlength{\parskip}{4pt}
\setlength{\parindent}{1.1em}
\setlist{itemsep=4pt,topsep=5pt}
\emergencystretch=2em

\begin{document}
% ed. (2026-09-29): no page anchors on the title page, which removes the
% duplicate hyperref destination page.1 (pdfTeX warning).
\hypersetup{pageanchor=false}
\begin{titlepage}
\thispagestyle{empty}
\vspace*{13mm}
{\color{MidnightBlue}\large TRANSSERIES \enspace / \enspace ASYMPTOTIC INVERSION}\par
\vspace{12mm}
{\Huge\bfseries Certified Inversion Through the\par $q\to1$ Transition}\par
\vspace{7mm}
{\Large Signed remainders, optimal truncation,\par and a modular cancellation window}\par
\vspace{8mm}
{\large With a shortest-block law for Gaussian multinomials}\par
\vspace{12mm}
\hrule
\vspace{5mm}
Prepared for Vladimir Reshetnikov\par
29 September 2026\par
Repository snapshot: \texttt{20301601562cf8}\par
\vspace{9mm}
\begin{abstract}
The simultaneous limit $q=\ee^h\to1^+$ and $x\to\infty$ reorganizes the
transseries of Gaussian binomial coefficients: the finite-product tail enters a
dilogarithmic dominant phase, while the Euler product contributes a separate
modular exponential. We prove an exact positive-kernel remainder formula for
the canonical real interpolation of the central Gaussian binomial. It yields
signed first-omitted-term bounds for every $h,x>0$, a fixed-order bound uniform
up to $hx=0$, and two-sided rational approximations from a determinate
Stieltjes measure.

The main quantitative result identifies the optimally truncated finite-size
remainder, including its first relative correction. Its leading value is
$(-1)^{M+1}\ee^{-2\pi x}/(\pi\sqrt{x})$ when $M=\pi x+O(1)$ and $hx$ stays
in a compact positive interval. Combining it with the modular term
$\ee^{-4\pi^2/h}$ gives a stable additive law for inverse errors. Their true
balance occurs in a logarithmically shifted window, not exactly at $hx=2\pi$;
even and odd truncations respectively cancel and reinforce. A multinomial
extension shows that the shortest block and its multiplicity determine the
leading optimal remainder. Complete proofs, reproducible exact-rational and
high-precision checks, a repository crosswalk, and a further-research agenda
are included.
\end{abstract}
\vfill
\begin{resultbox}
\small\textbf{Status and scope.} This is a proof-based research manuscript,
not an independently refereed or Lean-checked result. The explicit refinements
are developed relative to the inspected repository chapter. Classical
partial-fraction, modular, moment, and summation machinery is attributed;
worldwide priority for the refinements has not been established. No general
open conjecture about all transseries is claimed solved.
\end{resultbox}
\end{titlepage}
\hypersetup{pageanchor=true}

\tableofcontents
\newpage

\section{The problem, the contribution, and the boundary of novelty}
\label{sec:scope}

\subsection{A concrete gap rather than a generic request for a breakthrough}
The \repo\ companion \emph{Combinatorial Transseries and Their Inverses}
already contains a detailed chapter on the singular transition $q\to1$
\cite{repo}. Its forward expansion isolates a dilogarithmic phase, a logarithmic
amplitude, Bernoulli corrections, and an exact modular factor. It also gives
finite-order inverse corrections and an all-order formal inverse formula.
Those are starting points here, not new results of this manuscript.

Two distinctions in that chapter are particularly important. Its uniform
finite-order error theorem assumes $\tau=hx\ge\tau_0>0$, and it explicitly
leaves uniform matching to $\tau=0$ to a different analysis. It also states
that the modular factor is not the whole beyond-all-orders structure of the
finite-size tail. Consequently, retaining a modular term next to an unspecified
algebraic error does not establish which exponential actually controls an
inverse approximation. The present paper supplies a positive remainder kernel
and uses it to answer that more precise question.

The resulting advances are linked rather than a collection of unrelated
coefficient calculations. Positivity makes the divergent expansion enveloping.
The same positivity yields a Stieltjes measure and rational enclosures. A
large-moment analysis then determines the exact leading optimal remainder,
which permits a signed comparison with the modular correction. Finally,
inverse-error transport preserves that comparison because the interpolation
has a uniformly positive slope.

\subsection{What is proved here}
The principal results are the following. \Cref{thm:signed} gives an exact
signed remainder and a strict first-omitted-term bound. \Cref{cor:uniform}
extends fixed-order control to $0<hx\le\tau_1$ in the natural variable $x$.
\Cref{thm:stieltjes,thm:pade} give a positive Stieltjes representation and
convergent two-sided Pad\'e bounds. \Cref{thm:optimal} determines the optimal
remainder through relative order $x^{-1}$. \Cref{thm:inverse,thm:window}
give the two-exponential inverse-error law and its logarithmically shifted
cancellation window. \Cref{thm:shortest} extends the leading law to arbitrary
fixed positive multinomial block proportions.

These are assertions with proofs below. The numerical checks are separate:
they test signs, constants, rational bounds, and nonlinear residuals, but do not
replace the arguments. In particular, no unexecuted Lean code is presented as
formal verification.

\subsection{Relation to existing analysis}
The partial-fraction expansion of the hyperbolic cotangent and the modular
transformation of the Dedekind eta function are classical; convenient precise
references are DLMF~4.36.3 and~23.18.5 \cite{dlmf-hyper,dlmf-eta}. Modular
representations of general $q$-Pochhammer products involving the dilogarithm
and Laplace integrals were developed by Zhang \cite{zhang}. Modern resurgence
and summability results for weighted $q$-Pochhammer sums appear in Fantini and
Rella \cite{fantini-rella}. The existence of a Borel description of a
$q$-product is therefore not a novelty claim here.

Likewise, positivity, quadrature, moment determinacy, and Pad\'e convergence
belong to classical Stieltjes theory; Simon gives a substantial treatment of
its operator and approximation aspects \cite{simon}. Here the specific
measure, the resulting certificates for this interpolation, and their
connection to the inverse cancellation calculation are made explicit.

The sharp parity-sensitive inverse comparison is the main proposed research
contribution. A targeted comparison with the cited sources and the pinned
repository does not constitute an exhaustive priority search. The assertions
are deliberately formulated so that their correctness can be assessed without
accepting a claim that they are unprecedented.

% ed. (2026-09-29): editorial note added by ProveIt (sibling packages,
% a duplicate principal theorem, uncited repository results); not part of
% the delivered text.
\begin{quote}\small
\textbf{Editorial note (ProveIt, 2026-09-29).}
\emph{Sibling packages.} This article is one of five packages that answer
the same gap of the companion for the same interpolation and are to be merged
into its chapter on the singular transition $q\to1$ \cite{ed:siblings}: the
$q$-multinomial article (every positive ray; its interpolation at $r=2$,
$\boldsymbol a=(1,1)$ is \eqref{eq:H}, and the constant
$m_*/(2\pi\sqrt{\alpha x})$ of \cref{thm:shortest} is its sharp $h=0$
constant), the central Gaussian-binomial article, this article, and the
gamma-core and theta-resolved articles. The five write the interpolation in different normalizations ($q=\ee^h$ or $Q=\ee^{-h}$; remainders $R_K$, $r_K$, $R_M$ or $L_h-F_{h,K}$; smallest ray parameter $a_*$, $\alpha$ or $a_{\min}$); a normalization dictionary is left to the merge.

\emph{Duplicate principal theorem.} \Cref{thm:optimal,thm:window}
re-derive, independently, the principal theorem of the central
Gaussian-binomial package, which this article, written from an earlier
snapshot, does not cite. That article proves
$\min_Kr_K=\ee^{-2\pi x}/(\pi\sqrt x)\,(1+O_T(1/x))$, with every minimizing
index $\pi x+O_T(1)$, uniformly for $0\le hx\le T$: the leading term of
\eqref{eq:optimal} on a wider range. Its boundary layer
$hx=2\pi-(\log x)/(2x)+s'/x$, with optimal remainder over modular term
tending to $\ee^{-s'}/\pi$, is \eqref{eq:swindow} with
$s=s'+\log\pi+o(1)$, and its equal-size boundary agrees with
\eqref{eq:tauc}. The hypothesis here, $hx$ in a compact subset of
$(0,\infty)$, is narrower. These are one theorem with two independent
proofs; new here are the relative correction $C_1(d)/(2\pi x)$, the
Stieltjes--Pad\'e bounds, the parity-sensitive window and the exact
cancellation point.

\emph{Gaussian Coefficient Calculus.} Theorem \texttt{q3:thm:double-scaling}
of the Fabius tree's Gaussian-coefficient volume \cite{ed:gcc} expands
$\log\bigl[\begin{smallmatrix}n\\k\end{smallmatrix}\bigr]_{\ee^{-\tau/n}}$
to all orders only for $0<\tau_0\le\tau\le\tau_1$, the restriction of
\texttt{t3:prop:double-A}. Its central case $k/n=\frac12$ is this
interpolation at integers, with $n=2x$ and its $\tau$ equal to $2hx$;
\cref{cor:uniform} reaches $0<hx\le\tau_1$. This article does not cite that
volume.

\emph{Canonical volume and Lean.} \Cref{thm:optimal} is an explicit, sharp
instance of the canonical volume's general least-term theorem
\texttt{p0:thm:optimal-truncation} \cite{ed:tai}. The ``derivative lower
bound and root-existence bracket'' that \cref{sec:integration} asks for are,
in abstract form, the Lean theorems
\path{Fabius.exists_eq_in_residual_interval} and \path{Fabius.transport_bound}
\cite{ed:lean}; the special-function hypotheses are not formalized. This
article cites neither.
\end{quote}

\section{Exact interpolation and the two small mechanisms}
\label{sec:model}

For $0<Q<1$, write
\[
 \qpoch{z}{Q}=\prod_{j=0}^{\infty}(1-zQ^j).
\]
Throughout the central-binomial analysis,
\begin{equation}
 h>0,\qquad Q=\ee^{-h},\qquad q=\ee^h,\qquad x>0,
 \qquad \tau=hx,\qquad t=\ee^{-\tau}.
 \label{eq:parameters}
\end{equation}
All logarithms of positive quantities are real logarithms. Define
\begin{equation}
 P_h=\qpoch{Q}{Q},\qquad
 P_h(x)=\frac{\qpoch{Q}{Q}}{\qpoch{Q^{x+1}}{Q}},\qquad
 H_h(x)=\ee^{hx^2}\frac{P_h(2x)}{P_h(x)^2},\qquad
 L_h(x)=\log H_h(x).
 \label{eq:H}
\end{equation}
The distinction between the constant $P_h$ and the function $P_h(x)$ is
maintained throughout. At an integer $n\ge0$, $P_h(n)=\prod_{j=1}^n(1-Q^j)$,
and elementary reversal of the factors gives
\[
 H_h(n)=\genfrac{[}{]}{0pt}{}{2n}{n}_{q}.
\]
Thus \eqref{eq:H} specifies an interpolation, rather than assuming that an
integer sequence has a unique analytic continuation.

Introduce the exact tail
\begin{equation}
 A_h(t)=\sum_{m=1}^{\infty}\frac{t^m}{m(\ee^{hm}-1)}.
 \label{eq:A}
\end{equation}
Expanding $-\log(1-z)$ inside absolutely convergent products proves
\begin{equation}
 L_h(x)=hx^2-\log P_h+A_h(t^2)-2A_h(t).
 \label{eq:exactL}
\end{equation}
The modular identity supplies the other mechanism:
\begin{align}
 P_h&=\sqrt{\frac{2\pi}{h}}\,
       \exp\left(-\frac{\pi^2}{6h}+\frac h{24}\right)
       \qpoch{r_h}{r_h},&
 r_h&=\ee^{-4\pi^2/h},\label{eq:modular}\\
 \E(h)&=-\log\qpoch{r_h}{r_h}
        =\sum_{n=1}^{\infty}\sigma_{-1}(n)r_h^n,&
 \sigma_{-1}(n)&=\sum_{d\mid n}\frac1d.\label{eq:E}
\end{align}
For completeness, apply $\eta(-1/z)=\sqrt{-iz}\eta(z)$ at
$z=ih/(2\pi)$ and use $\eta(ih/(2\pi))=\ee^{-h/24}P_h$ to obtain
\eqref{eq:modular}. Expanding each logarithm in the residual Euler product
and grouping equal products of indices proves \eqref{eq:E}. In particular,
\begin{equation}
 \E(h)=r_h+O(r_h^2),\qquad
 0<\E(h)\le\frac{r_h}{(1-r_h)^2}.
 \label{eq:Ebound}
\end{equation}
The inequality follows from $\sigma_{-1}(n)\le n$.

Define the known phase and combined amplitude
\begin{align}
 S(\tau)&=\tau^2+\frac{\pi^2}{6}
          +\Li_2(\ee^{-2\tau})-2\Li_2(\ee^{-\tau}),\label{eq:S}\\
 G(h,\tau)&=\frac12\log\left(\frac{h}{2\pi}
                                  \coth\frac\tau2\right).\label{eq:G}
\end{align}
Differentiation gives useful elementary identities:
\begin{equation}
 S'(\tau)=2\log(1+\ee^\tau),\qquad
 S''(\tau)=\frac{2\ee^\tau}{1+\ee^\tau},\qquad S(0)=0.
 \label{eq:Sder}
\end{equation}
Hence $S$ maps $(0,\infty)$ increasingly onto $(0,\infty)$. These phase
identities and the formal inverse based on them already occur in the
repository \cite{repo}.

Let the Bernoulli numbers be specified by
$z/(\ee^z-1)=\sum_{n\ge0}B_nz^n/n!$ near zero. For an integer $M\ge0$ set
\begin{align}
 B_M(h,x)={}&\frac{S(\tau)}h+G(h,\tau)-\frac h{24}\nonumber\\
 &+\sum_{k=1}^{M}\frac{B_{2k}h^{2k-1}}{(2k)!}
       \left(\Li_{2-2k}(t^2)-2\Li_{2-2k}(t)\right).
 \label{eq:BM}
\end{align}
An empty sum means zero. Our remainder is defined by the exact equation
\begin{equation}
 L_h(x)=B_M(h,x)+\E(h)+R_M(h,x).
 \label{eq:Rdef}
\end{equation}
The notation $B_M$ denotes a truncation, not a Bernoulli number when supplied
with arguments $(h,x)$.

\section{A global signed remainder formula}
\label{sec:signed}

\begin{lemma}[Finite thermal-kernel remainder]
\label{lem:thermal}
For every real $z>0$ and integer $M\ge0$,
\begin{align}
 \frac1{\ee^z-1}
 &=\frac1z-\frac12+\sum_{j=1}^{M}\frac{B_{2j}z^{2j-1}}{(2j)!}
       +\rho_M(z),\label{eq:thermal}\\
 \rho_M(z)
 &=(-1)^M2z^{2M+1}\sum_{k=1}^{\infty}
   \frac{1}{(2\pi k)^{2M+2}\{1+(z/(2\pi k))^2\}}.
 \label{eq:rho}
\end{align}
\end{lemma}
\begin{proof}
The cotangent partial fraction, with its argument halved, yields
\[
 \frac1{\ee^z-1}=\frac1z-\frac12+
              2z\sum_{k\ge1}\frac1{(2\pi k)^2+z^2}.
\]
In each summand use the finite identity
\[
 \frac1{1+u}=\sum_{j=0}^{M-1}(-u)^j+\frac{(-u)^M}{1+u}.
\]
The coefficients agree with the Taylor coefficients at zero; equivalently,
\[\frac{B_{2j}}{(2j)!}=(-1)^{j-1}\frac{2\zeta(2j)}{(2\pi)^{2j}}.\]
The residual sum converges absolutely. This argument uses a \emph{finite}
geometric identity at arbitrary positive $z$, not an invalid globally
convergent Bernoulli expansion.
\end{proof}

\begin{theorem}[Exact signed envelope]
\label{thm:signed}
Let
\begin{equation}
 w_m(\tau)=2\ee^{-m\tau}-\ee^{-2m\tau}>0.
 \label{eq:w}
\end{equation}
For all $h,x>0$ and all integers $M\ge0$,
\begin{equation}
 \boxed{R_M(h,x)=(-1)^{M+1}2h^{2M+1}
 \sum_{k,m\ge1}\frac{w_m(\tau)m^{2M}}
 {(2\pi k)^{2M+2}\{1+(hm/(2\pi k))^2\}}.}
 \label{eq:exactR}
\end{equation}
In particular, with
\begin{equation}
 T_{M+1}(h,x)=\frac{|B_{2M+2}|h^{2M+1}}{(2M+2)!}
       \left(2\Li_{-2M}(t)-\Li_{-2M}(t^2)\right),
 \label{eq:T}
\end{equation}
one has the strict inequalities
\begin{equation}
 0<(-1)^{M+1}R_M(h,x)<T_{M+1}(h,x).
 \label{eq:envelope}
\end{equation}
\end{theorem}
\begin{proof}
Insert \eqref{eq:thermal} into \eqref{eq:A}. The terms $1/(hm)$ and
$-1/2$ sum to $\Li_2(t)/h$ and $\tfrac12\log(1-t)$.
The finite polynomial terms give the polylogarithms in \eqref{eq:BM}.
The remaining combination is
\[
 -\sum_{m\ge1}\frac{w_m(\tau)}m\rho_M(hm),
\]
which gives \eqref{eq:exactR}. All absolute sums are finite because
$w_m\le2\ee^{-m\tau}$ and the $k$-sum is bounded by a zeta value.
Every factor of the unsigned remainder is positive. Replacing
$\{1+(hm/(2\pi k))^2\}^{-1}$ by $1$ strictly increases every summand.
Summation and the Bernoulli--zeta identity give \eqref{eq:T}.
\end{proof}

For even $M$ the exact logarithm lies \emph{below} $B_M+\E$; for odd $M$ it
lies above. More explicitly,
\begin{align}
 B_M+\E-T_{M+1}<L_h<B_M+\E,
       &&M\text{ even},\label{eq:evenenv}\\
 B_M+\E<L_h<B_M+\E+T_{M+1},
       &&M\text{ odd}.\label{eq:oddenv}
\end{align}
This is stronger than an unspecified asymptotic $O$-term: it is valid before
taking any limit and records the sign needed later for cancellation.

\subsection{Uniform control at the classical endpoint}
\begin{corollary}[No positive lower cutoff on $hx$]
\label{cor:uniform}
For fixed $\tau_1>0$ and $M\ge0$, throughout $0<hx\le\tau_1$,
\begin{equation}
 |R_M(h,x)|\le
 \frac{4\zeta(2M+2)\ee^{\tau_1}(2M)!}
      {(2\pi)^{2M+2}x^{2M+1}}.
 \label{eq:uniform}
\end{equation}
Thus the same composite expansion has a uniform fixed-order remainder as
$x\to\infty$, even when $hx\to0$.
\end{corollary}
\begin{proof}
For every integer $p\ge0$ and $\tau>0$,
\[
 \sum_{m\ge1}m^p\ee^{-m\tau}
 \le\ee^\tau\int_0^\infty u^p\ee^{-\tau u}\dd u
 =\ee^\tau\frac{p!}{\tau^{p+1}}.
\]
Indeed on $[m,m+1]$, $u^p\ge m^p$ and
$\ee^{-\tau u}\ge\ee^{-\tau(m+1)}$. Apply this inequality with $p=2M$
to \eqref{eq:T}, use $w_m\le2\ee^{-m\tau}$ and $\tau=hx$.
\end{proof}

The endpoint is best understood with $S$ and $G$ retained intact. Their small
$\tau$ expansions are
\begin{align}
 S(\tau)&=2(\log2)\tau+\frac{\tau^2}{2}
                    +\frac{\tau^3}{12}+O(\tau^5),\label{eq:Ssmall}\\
 G(h,\tau)&=-\frac12\log(\pi x)+O(\tau^2).
 \label{eq:Gsmall}
\end{align}
For each fixed $k\ge1$,
\begin{equation}
 \frac{B_{2k}h^{2k-1}}{(2k)!}
 \left(\Li_{2-2k}(\ee^{-2\tau})-2\Li_{2-2k}(\ee^{-\tau})\right)
 \sim\frac{(2^{1-2k}-2)B_{2k}}
 {2k(2k-1)x^{2k-1}}
 \quad(\tau\downarrow0).
 \label{eq:classicalcoefs}
\end{equation}
These are the central-binomial Stirling coefficients. At $k=1$ the constant
part of the polylogarithm combination cancels the separately retained
$-h/24$. In fact their sum is
\begin{equation}
 -\frac h{24}+\frac h{12}\left(\Li_0(t^2)-2\Li_0(t)\right)
 =-h\frac{\ee^{2\tau}+4\ee^\tau+1}{24(\ee^{2\tau}-1)}
 =-\frac1{8x}-\frac{h\tau^3}{1440}+O(h\tau^5).
 \label{eq:firstcoefficient}
\end{equation}

\begin{remark}[What endpoint matching does not say]
The condition $hx\to0$ alone does not permit dropping all $q$-dependent
phase corrections. For example $\tau^2/(2h)=hx^2/2$ can be large even when
$hx$ is small. The uniform statement is \eqref{eq:uniform} for the composite
expression $B_M$, not an assertion that $H_h(x)$ is relatively close to the
classical binomial under only $hx\to0$. This distinction avoids a common
nonuniform interchange of limits.
\end{remark}

\section{Monotonicity and certified inverse transport}
\label{sec:inversecert}

\begin{proposition}[A slope bound independent of $h$]
\label{prop:slope}
The function $L_h$ is strictly increasing on $(0,\infty)$, extends
continuously with $L_h(0)=0$, and tends to infinity. For $x\ge1$,
\begin{equation}
 \mu:=2\log\frac32\ \le L_h'(x)\le S'(hx).
 \label{eq:slopebounds}
\end{equation}
\end{proposition}
\begin{proof}
Differentiate \eqref{eq:exactL}, justified by exponential convergence on
compact positive $x$-intervals:
\begin{equation}
 L_h'(x)=2hx+2h\sum_{m\ge1}
       \frac{\ee^{-hmx}-\ee^{-2hmx}}{\ee^{hm}-1}>0.
 \label{eq:exactderiv}
\end{equation}
The product definition gives continuity at zero; the $hx^2$ term in
\eqref{eq:exactL} gives divergence at infinity. The inequalities
\[
 \frac{\ee^{-hm}}m\le\frac h{\ee^{hm}-1}\le\frac1m
\]
follow from $1-\ee^{-z}\le z\le\ee^z-1$. Summing the lower estimate in
\eqref{eq:exactderiv} and combining the $2hx$ term with the logarithms gives
\[
 L_h'(x)\ge2\log\frac{\ee^{h(2x+1)}-1}{\ee^{h(x+1)}-1}
 \ge2\log\frac{2x+1}{x+1}.
\]
The second inequality follows because $(\ee^u-1)/u$ increases for $u>0$.
For $x\ge1$ the last ratio is at least $3/2$. The upper estimate sums to
$2hx+2\log(1+\ee^{-hx})=S'(hx)$.
\end{proof}

Given a positive target $Y$, put $y=\log Y$. The real inverse exists uniquely
when $Y\ge1$. For an inverse $x_*$ and a trial point $z$, both at least one,
\begin{equation}
 |z-x_*|\le\frac{|L_h(z)-y|}{\mu}.
 \label{eq:residual}
\end{equation}
This follows directly from the mean value theorem. With any rigorous
logarithmic enclosure $\ell_-(z)\le L_h(z)\le\ell_+(z)$, a usable bound is
\begin{equation}
 |z-x_*|\le\frac{\max\{|\ell_-(z)-y|,|\ell_+(z)-y|\}}{\mu}.
 \label{eq:intervaltransport}
\end{equation}
The assumption that the interval of interest lies in $[1,\infty)$ is explicit;
for smaller indices one uses a positive local slope minimum instead.

The signed envelopes \eqref{eq:evenenv}--\eqref{eq:oddenv} give immediate
choices of $\ell_\pm$. A bracket $u<v$ with $\ell_+(u)<y<\ell_-(v)$ contains
the unique inverse. This bracketing method does not require a truncated
expression itself to be globally monotone.

For the integer threshold $\min\{n:H_h(n)\ge Y\}$, the answer is
$\lceil x_*\rceil$. A certified bracket determines it when its endpoint
ceilings agree. If the bracket crosses an integer, evaluate the exact finite
Gaussian product there. Equality must not be decided from a rounded
floating-point approximation. These order-theoretic principles already
belong to the repository's inversion architecture \cite{repo}; the explicit
forward envelopes above are the additional input.

\section{A positive Stieltjes transform and rational enclosures}
\label{sec:stieltjes}

\begin{theorem}[The finite-size correction is Stieltjes]
\label{thm:stieltjes}
Fix $\tau>0$. Define, for $s\ge0$,
\begin{equation}
 \D_\tau(s)=2\sum_{k,m\ge1}\frac{w_m(\tau)}{(2\pi k)^2+s m^2}.
 \label{eq:D}
\end{equation}
Then
\begin{equation}
 L_h(\tau/h)=\frac{S(\tau)}h+G(h,\tau)-\frac h{24}
               +\E(h)-h\D_\tau(h^2).
 \label{eq:stieltjesL}
\end{equation}
More precisely,
\begin{align}
 \D_\tau(s)&=\int_{[0,\infty)}\frac{\dd\nu_\tau(u)}{1+su},\label{eq:measureD}\\
 \nu_\tau&=2\sum_{k,m\ge1}\frac{w_m(\tau)}{(2\pi k)^2}
                     \delta_{m^2/(2\pi k)^2},\label{eq:measure}\\
 \mu_j(\tau):=\int u^j\dd\nu_\tau(u)
 &=\frac{|B_{2j+2}|}{(2j+2)!}
       \left(2\Li_{-2j}(t)-\Li_{-2j}(t^2)\right).\label{eq:moments}
\end{align}
The measure is finite and positive, has infinite support, has all moments,
and is uniquely determined among measures on $[0,\infty)$ by those moments.
\end{theorem}
\begin{proof}
The case $M=0$ of \cref{lem:thermal} gives \eqref{eq:stieltjesL}. Factoring
$(2\pi k)^2$ from each denominator proves \eqref{eq:measureD}. Tonelli's
theorem evaluates the moments as a product of a zeta value and an exponential
power sum, giving \eqref{eq:moments}. Positivity is strict on infinitely many
support points, for example the points with $k=1$.

Here is a direct determinacy argument. For $0<c<2\pi\tau$,
\begin{equation}
 \int\ee^{c\sqrt u}\dd\nu_\tau(u)
 \le 4\sum_{k,m\ge1}\frac{\ee^{-m\tau+cm/(2\pi k)}}{(2\pi k)^2}<\infty.
 \label{eq:expsqrt}
\end{equation}
Also \cref{cor:uniform}'s moment estimate implies
$\mu_j\le C A^{2j}(2j)!$ for constants depending on $\tau$.
For any other positive measure with these moments, the monotone expansion of
$\cosh(c\sqrt u)$ gives a finite integral for sufficiently small $c$.
Its symmetric square-root lift to $\R$ therefore has an exponential moment
near zero. The moment generating function of that lift is analytic there
and determined by the moments; equality of the analytic transforms, followed
by uniqueness of the Fourier transform after exponential tilting, determines
the lift and hence the original measure. Equivalently, the Stieltjes
Carleman sum $\sum_j\mu_j^{-1/(2j)}$ diverges.
\end{proof}

For $s>0$, differentiation under the integral gives
\begin{equation}
 (-1)^j\D_\tau^{(j)}(s)
    =j!\int\frac{u^j\dd\nu_\tau(u)}{(1+su)^{j+1}}>0.
 \label{eq:complete}
\end{equation}
Thus $\D_\tau$ is completely monotone. All Hankel moment matrices are
positive definite: for any nonzero polynomial $p$,
$\int p(u)^2\dd\nu_\tau(u)>0$. The formal expansion is
\begin{equation}
 \D_\tau(s)\sim\sum_{j\ge0}(-1)^j\mu_j(\tau)s^j.
 \label{eq:formalD}
\end{equation}
It is generally divergent at every nonzero $s$ as a power series. Positivity
nevertheless gives convergent rational approximations.

\begin{theorem}[Two-sided Pad\'e certificates]
\label{thm:pade}
Let $[L/N]_\tau$ denote the Pad\'e approximant to \eqref{eq:formalD}, normalized
with denominator constant term one. For $N\ge1$ and $s>0$,
\begin{equation}
 [N-1/N]_\tau(s)<\D_\tau(s)<[N/N]_\tau(s).
 \label{eq:padebound}
\end{equation}
Both approximants exist, have no pole on $(0,\infty)$, and converge pointwise
to $\D_\tau(s)$ as $N\to\infty$. Consequently, replacing $\D_\tau$ in
\eqref{eq:stieltjesL} by the upper and lower rational bounds gives respectively
lower and upper bounds for $L_h$.
\end{theorem}
\begin{proof}
Take the $N$-node Gaussian quadrature for $\nu_\tau$. Its nodes are positive,
its weights are positive, and it is exact through degree $2N-1$. These facts
follow from the positive definite moment form and its degree-$N$ orthogonal
polynomial. Applied to $f_s(u)=(1+su)^{-1}$, it is a rational function with
numerator degree at most $N-1$ and denominator degree $N$, whose expansion
agrees through $s^{2N-1}$. It is therefore $[N-1/N]_\tau$.

The polynomial Hermite interpolant matching $f_s$ and $f_s'$ at the Gaussian
nodes lies below $f_s$ on $[0,\infty)$: its error has the sign of
$f_s^{(2N)}$ times a squared node polynomial, which is positive. Integrating
and using quadrature exactness proves the lower inequality. This argument is
valid on each finite interval containing the evaluation point; its global
integral is permitted because all moments exist.

For the upper bound use Gauss--Radau quadrature with the prescribed node zero
and $N$ further positive nodes. The latter are the zeros of the degree-$N$
orthogonal polynomial for $u\dd\nu_\tau(u)$. The quadrature is exact through
degree $2N$. The interpolant matches at zero and matches values and
derivatives at the other nodes. Its error has the sign of
$f_s^{(2N+1)}u$ times a squared polynomial, which is negative. Thus the
interpolant, and after integration the quadrature, is an upper bound. Its
rational expression is $[N/N]_\tau$. Infinite support makes both inequalities
strict.

For convergence, regard either sequence of quadratures as positive discrete
measures. Their total mass is $\mu_0$ and their first moment is eventually
$\mu_1$, so they are tight. Each higher moment is eventually exact; the next
higher exact moment supplies uniform integrability for passing that moment
to a weak limit. Every subsequential limit therefore has all moments
$\mu_j$. Determinacy in \cref{thm:stieltjes} forces the limit to be
$\nu_\tau$. The bounded continuous function $f_s$ then gives convergence of
the quadrature values. Positive nodes ensure the absence of positive poles.
\end{proof}

\begin{example}[Rational certificates at $t=1/2$]
With $\tau=\log2$,
\[
 \mu_0=\frac5{36},\qquad \mu_1=\frac{19}{1215},\qquad
 [0/1]_\tau(s)=\frac{5/36}{1+(76/675)s}.
\]
Every moment is rational, since negative-integer polylogarithms are rational
functions of $t$. Hence every finite Pad\'e bound can be built in exact
rational arithmetic. The computation in \cref{sec:computation} uses this fact;
it does not infer rigorous inequalities from decimal agreement.
\end{example}

\section{Borel reconstruction and the relevant singularity scale}
\label{sec:borel}

The measure also exposes the exponential scale without making a claim that
the entire modular completion is encoded in the algebraic coefficients.
Use the Borel convention
\[
 \mathcal B\!\left(\sum_{n\ge0}a_nh^{n+1}\right)(\xi)
   =\sum_{n\ge0}\frac{a_n}{n!}\xi^n.
\]
For the finite-size series $F_\tau(h)\sim-h\sum_{j\ge0}(-1)^j\mu_jh^{2j}$,
its Borel transform is initially
\begin{equation}
 \widehat F_\tau(\xi)
 =-2\sum_{k,m\ge1}\frac{w_m(\tau)}{(2\pi k)^2}
                       \cos\frac{m\xi}{2\pi k}.
 \label{eq:borel}
\end{equation}
This converges normally for $|\operatorname{Im}\xi|<2\pi\tau$.

\begin{proposition}[Real Borel sum and nearest singularities]
\label{prop:borel}
For $h>0$,
\begin{equation}
 \int_0^\infty\ee^{-\xi/h}\widehat F_\tau(\xi)\dd\xi
       =-h\D_\tau(h^2).
 \label{eq:borelsum}
\end{equation}
The transform admits a meromorphic continuation whose possible poles lie in
\begin{equation}
 \{4\pi^2k\ell\pm2\pi i k\tau:k\ge1,\ \ell\in\Z\}
 \ \cup
 \{4\pi^2k\ell\pm4\pi i k\tau:k\ge1,\ \ell\in\Z\}.
 \label{eq:poleset}
\end{equation}
The nearest singularities to the origin are the nonzero simple poles
$\xi=\pm2\pi i\tau$.
\end{proposition}
\begin{proof}
On the real $\xi$-axis, the absolute value of the double series is bounded by
$2\sum_{k,m}w_m/(2\pi k)^2<\infty$. Fubini and
\[
 \int_0^\infty\ee^{-\xi/h}\cos(b\xi)\dd\xi
       =\frac h{1+h^2b^2}
\]
prove \eqref{eq:borelsum}. For continuation write
\[
 C_r(z)=\frac12\left(\frac{r\ee^{iz}}{1-r\ee^{iz}}
                       +\frac{r\ee^{-iz}}{1-r\ee^{-iz}}\right).
\]
Then
\[
 \widehat F_\tau(\xi)=-2\sum_{k\ge1}(2\pi k)^{-2}
      \left\{2C_t(\xi/(2\pi k))-C_{t^2}(\xi/(2\pi k))\right\}.
\]
For sufficiently large $k$, the argument lies uniformly in a small
neighborhood of zero on any fixed compact set; the remaining $k^{-2}$
series converges normally. The finite set of small-$k$ summands gives the
listed possible poles. At $\pm2\pi i\tau$, only the $k=1$, $C_t$ term has a
pole, so cancellation is impossible. No listed pole has smaller modulus.
\end{proof}

Some farther poles in \eqref{eq:poleset} can cancel, so that set is not
asserted to be an exact divisor. The distance $2\pi\tau$ from the origin
predicts the magnitude $\ee^{-2\pi\tau/h}=\ee^{-2\pi x}$ of a least-term
error. It does not imply that a new independent positive-real transseries
parameter of that size occurs in the exact function: on the positive axis
\eqref{eq:borelsum} already defines the finite-size contribution unambiguously.
The optimal error is the difference between that sum and a finite truncation.
The modular term $\E(h)$ is a separate, exactly specified flat contribution.

This distinction between a truncation error, a Borel singularity, and a
separately normalized nonperturbative term is essential. The cancellation
studied below is a real approximation-error phenomenon, not by itself a
Stokes jump.

\section{Sharp optimal truncation, with a first correction}
\label{sec:optimal}

In this section $h\downarrow0$, $\tau=hx$ belongs to a fixed compact set
$K\subset(0,\infty)$, and the integer $M$ satisfies $|M-\pi x|\le C$ for a
fixed constant $C$. Define
\begin{equation}
 a=2\pi x,\qquad d=2M-a,\qquad
 C_1(d)=\frac{d^2}{2}-\frac d2-\frac5{12},\qquad
 a_x=\frac{\ee^{-2\pi x}}{\pi\sqrt{x}}.
 \label{eq:ad}
\end{equation}
Here $a_x$ is an error scale; it is unrelated to the block lengths introduced
later.

\begin{lemma}[Uniform large exponential moments]
\label{lem:moments}
For integers $p\ge1$, let $Z_p(\tau)=\sum_{m\ge1}m^p\ee^{-m\tau}$. Then
\begin{equation}
 Z_p(\tau)=p!\sum_{\ell\in\Z}(\tau+2\pi i\ell)^{-p-1}.
 \label{eq:Zexact}
\end{equation}
For $\tau\in K$, there is $0<\rho_K<1$ such that, as $p\to\infty$,
\begin{equation}
 Z_p(\tau)=\frac{p!}{\tau^{p+1}}
                       \left(1+O_K(\rho_K^p)\right).
 \label{eq:Zasym}
\end{equation}
The same assertion holds for $Z_{p+j}/Z_p$ with each fixed nonnegative $j$,
after dividing by the corresponding factorial and power ratio.
\end{lemma}
\begin{proof}
The bilateral partial fraction for the derivative of the thermal kernel is
\[
 \frac{\ee^{-\tau}}{(1-\ee^{-\tau})^2}
       =\sum_{\ell\in\Z}(\tau+2\pi i\ell)^{-2}.
\]
It is also the $\operatorname{csch}^2$ partial fraction. Differentiating
$p-1$ times gives \eqref{eq:Zexact}; all derivative series converge normally
on compact positive $\tau$-intervals. The $\ell=0$ term is the displayed
leading term. For the rest, the ratios
$|\tau/(\tau+2\pi i\ell)|$ are uniformly below one for $\ell\ne0$.
Bounding two powers by a summable $\ell^{-2}$ sequence and the remaining
powers by a constant less than one proves \eqref{eq:Zasym}. The same estimate
with $p+j$ proves the final assertion.
\end{proof}

\begin{theorem}[Leading optimal remainder and its relative correction]
\label{thm:optimal}
Under the compactness and bounded-offset hypotheses above,
\begin{equation}
 \boxed{R_M(h,x)=(-1)^{M+1}a_x
       \left(1+\frac{C_1(d)}{2\pi x}+O_{K,C}(x^{-2})\right).}
 \label{eq:optimal}
\end{equation}
The first omitted term satisfies
\begin{equation}
 T_{M+1}(h,x)=2a_x
 \left(1+\frac{d^2/2+d/2+1/12}{2\pi x}
                         +O_{K,C}(x^{-2})\right).
 \label{eq:optimalT}
\end{equation}
In particular $|R_M|/T_{M+1}\to1/2$, uniformly in this regime.
\end{theorem}
\begin{proof}
Put $p=2M=a+d$. In \eqref{eq:exactR}, the contribution with $k=1$ and
$w_m$ replaced by $2\ee^{-m\tau}$ is
\begin{equation}
 V=\frac{4h^{p+1}}{(2\pi)^{p+2}}Z_p(\tau)
     \mathbb E_p f(Y),\qquad
 f(v)=\frac1{1+v^2},\quad Y=\frac{hm}{2\pi},
 \label{eq:V}
\end{equation}
where $\mathbb P_p(m)=m^p\ee^{-m\tau}/Z_p(\tau)$.

First control the discarded contributions. Replacing $\ee^{-m\tau}$ by
$\ee^{-2m\tau}$ reduces the leading large moment by $2^{-p-1}$, up to a
uniform relative error exponentially small in $p$. The sum over $k\ge2$ is
bounded by its unattenuated moment times
$\sum_{k\ge2}k^{-p-2}=O(2^{-p})$. The main expectation tends to $1/2$, as
shown next, so both discarded parts are exponentially small relative to $V$.
This also justifies keeping a sign when computing the main contribution.

The moment-ratio statement in \cref{lem:moments} gives
\begin{align*}
 \mathbb E_p(Y-1)&=\frac{d+1}{a}+O_{K,C}(\rho^p),\\
 \mathbb E_p(Y-1)^2&=\frac1a+O_{K,C}(a^{-2}),\\
 \mathbb E_p(Y-1)^3&=O_{K,C}(a^{-2}),\qquad
 \mathbb E_p(Y-1)^4=O_{K,C}(a^{-2}).
\end{align*}
Polynomial factors multiplying the exponential errors may be absorbed by
increasing $\rho<1$. All derivatives of $f$ needed here are bounded on
$[0,\infty)$. Taylor expansion through degree three at one, with the fourth
moment bounding the remainder, therefore yields
\begin{equation}
 \mathbb E_p f(Y)=\frac12-
rac{d/2+1/4}{a}
                         +O_{K,C}(a^{-2}),
 \label{eq:Ef}
\end{equation}
since $f'(1)=-1/2$ and $f''(1)=1/2$.

Using \eqref{eq:Zasym} in the factor before the expectation gives
\[
 \frac{4h^{p+1}}{(2\pi)^{p+2}}Z_p(\tau)
 =\frac2\pi\frac{\Gamma(p+1)}{a^{p+1}}
                  \left(1+O_K(\rho^p)\right).
\]
Stirling's formula at $p=a+d$, uniformly for bounded $d$, reads
\begin{equation}
 \frac{\Gamma(a+d+1)}{a^{a+d+1}}
 =\frac{\ee^{-a}}{\sqrt{x}}
   \left(1+\frac{d^2/2+d/2+1/12}{a}+O_C(a^{-2})\right).
 \label{eq:stirlingshift}
\end{equation}
Multiplying \eqref{eq:Ef} by this expansion gives precisely
$C_1(d)=d^2/2-d/2-5/12$ and proves \eqref{eq:optimal}. Without the factor
$f(Y)$, the same calculation gives \eqref{eq:optimalT}. All omitted
exponential errors are smaller than $a^{-2}$ uniformly on $K$.
\end{proof}

\begin{remark}[Choosing the index and choosing a parity]
The theorem permits the nearest integer to $\pi x$, or the nearest even or
odd integer. An offset bounded independently of $x$ changes only the
relative correction, not the leading exponential or its prefactor. The
quadratic polynomial in \eqref{eq:optimal} also distinguishes minimizing the
actual remainder from minimizing just the first omitted term. At the
continuous-offset level their first correction has its minimum at $d=1/2$
and $d=-1/2$, respectively. This observation does not assert a global discrete
minimum for every finite $h,x$.
\end{remark}

\begin{remark}[A regime not covered by this theorem]
Fixed $h$ with $x\to\infty$ forces $\tau\to\infty$ and is not covered by the
compactness hypothesis. The power-sum saddle then need not be approximated
uniformly by its continuous gamma analogue. The fixed-order signed bounds
remain valid, but the sharp prefactor in \eqref{eq:optimal} is not asserted in
that different limiting regime.
\end{remark}

\section{The two-exponential law for inverse error}
\label{sec:competition}

\subsection{Local inverse approximants}
Fix an exact inverse point by setting $y=L_h(x)$. The truncation order is
held fixed while the approximate equation is solved. Let $\widehat x_M$ solve
\begin{equation}
 B_M(h,\widehat x_M)=y
 \label{eq:uncorrectedroot}
\end{equation}
near $x$; this approximation omits the modular completion. Let
$\widetilde x_M$ solve the corrected equation
\begin{equation}
 B_M(h,\widetilde x_M)+\E(h)=y.
 \label{eq:correctedroot}
\end{equation}
These are local roots; no global invertibility of every truncation is needed.
The next theorem proves their existence and uniqueness in a neighborhood
containing the relevant exponentially small displacement.

\begin{lemma}[Derivative control near optimal truncation]
\label{lem:dercontrol}
Under the hypotheses of \cref{thm:optimal}, uniformly on a sufficiently small
fixed neighborhood of $x$,
\begin{equation}
 L_h'(z)=S'(hz)+O_K(h),\qquad
 |\partial_z^jR_M(h,z)|=O_{K,C}(a_x)\quad(j=0,1,2),
 \label{eq:dercontrol}
\end{equation}
and $\partial_z B_M(h,z)=S'(hz)+O_{K,C}(h)>0$ for sufficiently small $h$.
\end{lemma}
\begin{proof}
Differentiate \eqref{eq:stieltjesL} along $\tau=hz$. The derivative of its
amplitude is $O_K(h)$; the derivative of $-h\D_\tau(h^2)$ is $O_K(h^2)$,
by exponential domination of each differentiated weight. This proves the
first assertion. In \eqref{eq:exactR}, $z$ enters only through $w_m(hz)$.
Each derivative adds at most a fixed power of $hm$, and the derivatives of
$w_m$ are bounded by fixed multiples of $h^jm^j\ee^{-hmz}$ on the stated
neighborhood. The proof of \cref{thm:optimal}, with $p$ replaced by $p+j$,
bounds the resulting moments by a constant times $a_x$. Their denominator
factors are positive and at most one. Finally differentiate \eqref{eq:Rdef}.
The modular term is independent of $z$, and $S'$ has a positive compact
minimum.
\end{proof}

\begin{theorem}[Stable additive inverse-error law]
\label{thm:inverse}
Let $\tau=hx\in K\Subset(0,\infty)$ and $|M-\pi x|\le C$ as $h\downarrow0$.
The local roots \eqref{eq:uncorrectedroot} and \eqref{eq:correctedroot} exist
and are unique in a fixed small neighborhood of $x$. With
$\sigma_M=(-1)^{M+1}$,
\begin{align}
 \boxed{\widehat x_M-x
 =\frac{r_h+\sigma_M a_x}{S'(\tau)}
                         +O_{K,C}\bigl(h(r_h+a_x)\bigr),}
 \label{eq:inverseerror}\\
 \widetilde x_M-x
 &=\frac{\sigma_M a_x}{S'(\tau)}(1+O_{K,C}(h)).
 \label{eq:correctederror}
\end{align}
The first formula is an \emph{additive} estimate, valid even when its two
leading contributions nearly cancel.
\end{theorem}
\begin{proof}
At $x$, the residual for \eqref{eq:uncorrectedroot} is
$B_M(h,x)-y=-\E(h)-R_M(h,x)$. By \cref{lem:dercontrol} the derivative is
bounded below on a small fixed interval. Evaluating at points at distance
$C_0(r_h+a_x)$ on each side, with $C_0$ sufficiently large, brackets a zero;
strict positivity of the derivative proves uniqueness. The mean value
theorem then gives
\[
 \widehat x_M-x=\frac{\E(h)+R_M(h,x)}
                         {\partial_zB_M(h,\xi)}
\]
for an intermediate $\xi$. Since $\xi-x=O(r_h+a_x)$, the denominator is
$S'(\tau)+O(h)$, uniformly. Insert \eqref{eq:Ebound} and
\eqref{eq:optimal}. Here $x^{-1}=O_K(h)$ and $r_h^2$ is absorbed by $hr_h$.
This proves \eqref{eq:inverseerror}, with an error proportional to the sum
of the magnitudes rather than to their possibly vanishing difference.
For the corrected equation the residual is only $-R_M$ and the same proof
applies.
\end{proof}

If $\tau$ remains in a compact subinterval of $(0,2\pi)$, then
$a_x/r_h\to\infty$, so optimal finite-size truncation error dominates.
If it remains in a compact subinterval of $(2\pi,\infty)$, the modular term
dominates. At the nominal crossing $\tau=2\pi$, however,
\begin{equation}
 \frac{a_x}{r_h}=\frac1{\pi\sqrt{x}}
          =\frac{\sqrt h}{\pi\sqrt{2\pi}}\longrightarrow0.
 \label{eq:nominalcross}
\end{equation}
Equal exponential actions do not mean equal errors; their algebraic
prefactors matter.

\subsection{The logarithmically shifted window}
Define the exact balance coordinate
\begin{equation}
 s=s(h,x):=2\pi x-\frac{4\pi^2}{h}
                   +\frac12\log x+\log\pi.
 \label{eq:swindow}
\end{equation}
It is chosen so that
\begin{equation}
 \frac{a_x}{r_h}=\ee^{-s}
 \label{eq:balanceratio}
\end{equation}
holds identically, not merely asymptotically.

\begin{theorem}[Parity-sensitive cancellation window]
\label{thm:window}
Suppose $s$ stays in a fixed bounded interval and $|M-\pi x|\le C$.
Then
\begin{equation}
 \boxed{\frac{S'(hx)}{r_h}(\widehat x_M-x)
         =1+(-1)^{M+1}\ee^{-s}+O_C(h).}
 \label{eq:profile}
\end{equation}
For even $M$, the leading profile is $1-\ee^{-s}$; for odd $M$, it is
$1+\ee^{-s}$. Let $x_c(h)$ be the unique positive solution of $s(h,x)=0$.
Then
\begin{align}
 x_c(h)&=\frac{2\pi}{h}+\frac{\log h}{4\pi}
            -\frac{\log(\pi\sqrt{2\pi})}{2\pi}
              +O(h|\log h|),\label{eq:xc}\\
 hx_c(h)&=2\pi+\frac{h\log h}{4\pi}
            -\frac{h\log(\pi\sqrt{2\pi})}{2\pi}
              +O(h^2|\log h|).\label{eq:tauc}
\end{align}
At $x=x_c(h)$ and even $M$,
\begin{equation}
 \frac{S'(hx)}{r_h}(\widehat x_M-x)
       =-\frac{C_1(2M-2\pi x)}{2\pi x}+O_C(h^2).
 \label{eq:refinedwindow}
\end{equation}
In particular the uncorrected inverse error is $O(hr_h)$ there.
\end{theorem}
\begin{proof}
Bounded $s$ implies $hx\to2\pi$, as follows immediately from
\eqref{eq:swindow}; in particular the compactness assumptions of
\cref{thm:inverse} hold. Divide \eqref{eq:inverseerror} by $r_h/S'(hx)$ and
use \eqref{eq:balanceratio}. The function $s(h,x)$ strictly increases from
$-\infty$ to $+\infty$ in $x$, giving the unique center.

For its expansion write $x=2\pi/h+u$. The defining equation becomes
\[
 2\pi u+\frac12\log(2\pi/h)+\log\pi
         +\frac12\log\left(1+\frac{hu}{2\pi}\right)=0.
\]
It first gives $u=O(|\log h|)$ and then \eqref{eq:xc}; multiplication by
$h$ proves \eqref{eq:tauc}.

At $s=0$, $a_x=r_h$. For even $M$, \cref{thm:optimal} and
\eqref{eq:Ebound} give
\[
 \E(h)+R_M(h,x)=-r_h\frac{C_1(d)}{2\pi x}+O_C(h^2r_h).
\]
The root displacement is already $O(hr_h)$. Dividing its residual by
$S'(hx)+O(h)$ changes the normalized result only by $O(h^2)$, which proves
\eqref{eq:refinedwindow}.
\end{proof}

The gain $O(r_h)\to O(hr_h)$ is an algebraic gain at the same exponential
action. It is not a proof that an entire exponential sector disappears.
The dependence on $2M-2\pi x$ also explains why the refined error oscillates
with integer truncation choices instead of varying monotonically in $h$.

\subsection{An exact cancellation point, not only a formal one}
\begin{proposition}[Unique exact cancellation for fixed even order]
\label{prop:exactcancel}
For every fixed $h>0$ and every even integer $M\ge2$, there is a unique
$x_{h,M}>0$ such that
\begin{equation}
 B_M(h,x_{h,M})=L_h(x_{h,M}).
 \label{eq:exactcancel}
\end{equation}
Choose even $M=M(h)$ with $|M-\pi x_c(h)|\le C$. Then, with
$d_c=2M-2\pi x_c(h)$,
\begin{equation}
 x_{h,M}=x_c(h)+\frac{C_1(d_c)}{4\pi^2x_c(h)}+O_C(h^2).
 \label{eq:exactcancellocation}
\end{equation}
\end{proposition}
\begin{proof}
Let $U_M=(-1)^{M+1}R_M>0$. For fixed $h,M$, every $w_m(hx)$ strictly
 decreases in $x$, since
$\partial_xw_m=-2hm(\ee^{-hmx}-\ee^{-2hmx})<0$.
Therefore \eqref{eq:exactR} makes $U_M$ continuous and strictly decreasing.
It tends to zero at infinity. For $M\ge1$, its $k=1$ part tends to infinity
at zero: for large $m$ the polynomial factor
$m^{2M}/(1+(hm/(2\pi))^2)$ is bounded below by a positive multiple of
$m^{2M-2}$, and $w_m(hx)\to1$ for each $m$. Monotone convergence then gives
divergence. Thus $U_M=\E(h)$ has exactly one solution. For even $M$ this is
\eqref{eq:exactcancel}.

Under the bounded-offset choice, the window theorem and monotonicity first
put the solution at $s=O(h)$. Use \eqref{eq:optimal} there to get
\[
 s(h,x_{h,M})=\log\left(1+\frac{C_1(2M-2\pi x_{h,M})}{2\pi x_{h,M}}
                                  +O_C(h^2)\right)
               =\frac{C_1(d_c)}{2\pi x_c}+O_C(h^2).
\]
Here moving $x$ by $O(h)$ changes the displayed coefficient by $O(h^2)$.
Since $\partial_xs=2\pi+1/(2x)$, expansion around $x_c$ proves
\eqref{eq:exactcancellocation}.
\end{proof}

There is no corresponding zero for an odd truncation: then
$\E+R_M>0$ everywhere. The exact cancellation point concerns the logarithmic
forward residual and consequently any local inverse built at that target.
It does not require the truncated expression to be globally invertible.

\section{Gaussian multinomials and the shortest-block law}
\label{sec:multi}

Let $b\ge2$ be the number of blocks, let
$\boldsymbol\alpha=(\alpha_1,\ldots,\alpha_b)$ be fixed positive real
numbers, and put
\begin{equation}
 A=\sum_{i=1}^b\alpha_i,\quad
 c=\frac12\left(A^2-\sum_{i=1}^b\alpha_i^2\right),\quad
 \alpha=\min_i\alpha_i,\quad m_* =\#\{i:\alpha_i=\alpha\}.
 \label{eq:alphas}
\end{equation}
Define
\begin{equation}
 H_{h,\boldsymbol\alpha}(x)=\ee^{hc x^2}
                   \frac{P_h(Ax)}{\prod_iP_h(\alpha_i x)},\qquad
 L_{h,\boldsymbol\alpha}(x)=\log H_{h,\boldsymbol\alpha}(x).
 \label{eq:multidef}
\end{equation}
When the block sizes $\alpha_i x$ are integers, this is the usual Gaussian
multinomial at $q=\ee^h$. Real positive block lengths specify its interpolation.
The central case is $b=2$ and $(\alpha_1,\alpha_2)=(1,1)$.

The exact phase and amplitude are
\begin{align}
 S_{\boldsymbol\alpha}(\tau)
 &=c\tau^2+(b-1)\frac{\pi^2}{6}
       +\Li_2(\ee^{-A\tau})-\sum_i\Li_2(\ee^{-\alpha_i\tau}),
       \label{eq:multiS}\\
 G_{\boldsymbol\alpha}(h,\tau)
 &=\frac12\log\left\{
  \left(\frac h{2\pi}\right)^{b-1}
   \frac{1-\ee^{-A\tau}}{\prod_i(1-\ee^{-\alpha_i\tau})}\right\}.
       \label{eq:multiG}
\end{align}
Replace the corresponding phase, amplitude, and polylogarithm combination in
\eqref{eq:BM} to define
\begin{align}
 B_M^{\boldsymbol\alpha}(h,x)
 ={}&\frac{S_{\boldsymbol\alpha}(\tau)}h+
 G_{\boldsymbol\alpha}(h,\tau)-\frac{(b-1)h}{24}\nonumber\\
 &+\sum_{k=1}^{M}\frac{B_{2k}h^{2k-1}}{(2k)!}
 \left\{\Li_{2-2k}(\ee^{-A\tau})-
                        \sum_i\Li_{2-2k}(\ee^{-\alpha_i\tau})\right\}.
 \label{eq:multiB}
\end{align}
Define $R_M^{\boldsymbol\alpha}$ by
\begin{equation}
 L_{h,\boldsymbol\alpha}(x)=B_M^{\boldsymbol\alpha}(h,x)
                    +(b-1)\E(h)+R_M^{\boldsymbol\alpha}(h,x).
 \label{eq:multiRdef}
\end{equation}

\begin{proposition}[Positive kernel and uniform inverse slope]
\label{prop:multikernel}
All signed-envelope and Stieltjes conclusions above hold for the multinomial
on replacing $w_m$ by
\begin{equation}
 w_m^{\boldsymbol\alpha}(\tau)=\sum_i\ee^{-\alpha_i\tau m}
                                   -\ee^{-A\tau m}>0.
 \label{eq:multiw}
\end{equation}
The phase slope is
\begin{equation}
 S_{\boldsymbol\alpha}'(\tau)=
 \sum_i\alpha_i\log\frac{\ee^{A\tau}-1}{\ee^{\alpha_i\tau}-1}>0.
 \label{eq:multislope}
\end{equation}
For $x\ge1/\alpha$,
\begin{equation}
 0<\mu_{\boldsymbol\alpha}:=
   \sum_i\alpha_i\log\frac{A+\alpha}{\alpha_i+\alpha}
 \le L_{h,\boldsymbol\alpha}'(x)
 \le S_{\boldsymbol\alpha}'(hx).
 \label{eq:multislopebound}
\end{equation}
\end{proposition}
\begin{proof}
The exact logarithm is
$hc x^2+(1-b)\log P_h+A_h(\ee^{-A\tau})-
\sum_iA_h(\ee^{-\alpha_i\tau})$. Inserting the thermal remainder gives
exactly \eqref{eq:exactR} with the new weight. It is positive, and is bounded
above by $b\ee^{-\alpha\tau m}$; every convergence, moment, and quadrature
argument therefore applies. The modular multiplicity is $b-1$.

Differentiating the phase and using $A=\sum_i\alpha_i$ and
$2c=A^2-\sum_i\alpha_i^2$ proves \eqref{eq:multislope}. For the exact
logarithm the positive derivative correction is
\[
 h\sum_{m\ge1}\sum_i\alpha_i
  \frac{\ee^{-\alpha_i hxm}-\ee^{-Ahxm}}{\ee^{hm}-1}.
\]
The same lower thermal bound as in \cref{prop:slope}, combined with the
quadratic derivative, yields
\[
 L_{h,\boldsymbol\alpha}'(x)
 \ge\sum_i\alpha_i\log
       \frac{\ee^{h(Ax+1)}-1}{\ee^{h(\alpha_i x+1)}-1}
 \ge\sum_i\alpha_i\log\frac{Ax+1}{\alpha_i x+1}.
\]
Every last ratio increases with $x$. Evaluating at $x=1/\alpha$ gives the
lower bound. The upper thermal bound gives the phase slope.
\end{proof}

\begin{theorem}[Shortest-block universality]
\label{thm:shortest}
Fix $\boldsymbol\alpha$ and $K\Subset(0,\infty)$. Suppose $h\downarrow0$,
$hx\in K$, and $|M-\pi\alpha x|\le C$. Then
\begin{equation}
 \boxed{R_M^{\boldsymbol\alpha}(h,x)
 =(-1)^{M+1}\frac{m_*}{2\pi\sqrt{\alpha x}}\ee^{-2\pi\alpha x}
                         (1+O_{K,C,\boldsymbol\alpha}(h)).}
 \label{eq:shortest}
\end{equation}
For a target $y=L_{h,\boldsymbol\alpha}(x)$, the local root of
$B_M^{\boldsymbol\alpha}(h,\widehat x)=y$ satisfies
\begin{align}
 \widehat x-x
 ={}&\frac{(b-1)r_h+(-1)^{M+1}a_{\boldsymbol\alpha,x}}
             {S_{\boldsymbol\alpha}'(hx)}\nonumber\\
 &+O_{K,C,\boldsymbol\alpha}\left(h\{r_h+a_{\boldsymbol\alpha,x}\}\right),
 \qquad
 a_{\boldsymbol\alpha,x}:=
        \frac{m_*\ee^{-2\pi\alpha x}}{2\pi\sqrt{\alpha x}}.
 \label{eq:multiinverse}
\end{align}
The normalized cancellation coordinate is
\begin{equation}
 s_{\boldsymbol\alpha}
 =2\pi\alpha x-\frac{4\pi^2}{h}+\frac12\log(\alpha x)
                       +\log\frac{2\pi(b-1)}{m_*}.
 \label{eq:multis}
\end{equation}
On bounded $s_{\boldsymbol\alpha}$-intervals the profile is again
$1+(-1)^{M+1}\ee^{-s_{\boldsymbol\alpha}}$, after normalizing the inverse
error by $S_{\boldsymbol\alpha}'(hx)/((b-1)r_h)$.
\end{theorem}
\begin{proof}
In the power moment of \eqref{eq:multiw}, the $m_*$ shortest blocks contribute
$m_* Z_p(\alpha\tau)$. Every other block contributes a relative factor
$(\alpha/\alpha_i)^{p+1}$, up to an exponentially small relative error;
the subtracted total block contributes $(\alpha/A)^{p+1}$. The fixed positive
gaps make these negligible. The $k\ge2$ remainder is again exponentially
small in $p$.

Apply the calculation \eqref{eq:V}--\eqref{eq:Ef} with
$a=2\pi\alpha x$ and the leading weight multiplicity $m_*$ instead of two.
The prefactor before its limiting expectation $1/2$ is
$(m_*/\pi)\Gamma(p+1)/a^{p+1}$, giving \eqref{eq:shortest}.
The same differentiated-moment estimates and \cref{prop:multikernel}
justify local inverse transport exactly as in \cref{thm:inverse}.
Finally $a_{\boldsymbol\alpha,x}/((b-1)r_h)=\ee^{-s_{\boldsymbol\alpha}}$
is an identity, proving the profile.
\end{proof}

This law says more than that a multinomial has a similar kind of expansion.
At optimal order, the longest and intermediate blocks disappear from the
leading finite-size error; the shortest block controls the exponential action,
and its multiplicity controls the prefactor. The remaining blocks still
influence the phase slope and the modular multiplicity. Uniformity as two
shortest lengths coalesce is a different, explicitly open problem below.

\section{Reproducible checks and an implementation route}
\label{sec:computation}

\subsection{What the supplied program actually verifies}
The package includes the script \texttt{verify.py} and its recorded output,
\texttt{results.json}, in the \texttt{verification} directory. The script uses exact rational
arithmetic for Pad\'e coefficients and Hankel determinants, and variable
high-precision arithmetic for the analytic checks. It checks positive
remainders against the direct thermal series, compares optimal remainder
constants, checks the parity-even inverse window with a nonlinear residual,
and checks three multinomial examples. All included assertions passed in the
recorded run.

The analytic computations use \texttt{mpmath}, not outward-rounded interval
arithmetic. Thus the output is evidence of reproducibility and error
localization, not a certified numerical proof of every reported decimal.
The inequalities are justified by the preceding theorems. A production
certifier must combine them with interval evaluations of elementary functions,
polylogarithms, and the modular product.

\subsection{Signed remainder checks}
\begin{table}[htbp]
\centering
\begin{tabular}{@{}rrrrr@{}}
\toprule
$h$ & $x$ & $M$ & $|R_M|$ & $|R_M|/T_{M+1}$\\
\midrule
$0.4$ & $2.5$ & $0$ & $3.3262523812\cdot10^{-2}$ & $0.9905105051$\\
$0.4$ & $2.5$ & $1$ & $3.1866855465\cdot10^{-4}$ & $0.9567940861$\\
$0.1$ & $10$ & $3$ & $1.1693752480\cdot10^{-10}$ & $0.9861299427$\\
$0.001$ & $10$ & $4$ & $1.6447426745\cdot10^{-12}$ & $0.9779321605$\\
$1$ & $3$ & $6$ & $5.4858162652\cdot10^{-9}$ & $0.6822534171$\\
\bottomrule
\end{tabular}
\caption{Direct thermal-series remainders obey the signed first-omitted-term
bound. The fourth row probes $hx=0.01$, outside any fixed lower cutoff such
as $hx\ge0.1$.}
\label{tab:signed}
\end{table}

The direct reference value sums
\begin{equation}
 -\sum_{m\ge1}\frac{w_m(\tau)}m
 \left\{\frac1{\ee^{hm}-1}-\frac1{hm}+\frac12\right\},
 \label{eq:thermalreference}
\end{equation}
not the Bernoulli expansion being tested. Subtraction from the truncated
expression checks both the residual sign and its magnitude. A truncated
positive double sum is separately compared using the analytic tail bounds in
\cref{app:tails}.

\subsection{The sharp prefactor and half-next-term rule}
\begin{table}[htbp]
\centering
\begin{tabular}{@{}rrrrrr@{}}
\toprule
$\tau$ & $h$ & $M$ & $|R_M|/a_x$ & $|R_M|/T_{M+1}$ & $1+C_1(d)/(2\pi x)$\\
\midrule
$1$ & $0.5$ & $6$ & $0.9999174546$ & $0.5015392405$ & $1.0021411494$\\
$1$ & $0.25$ & $13$ & $0.9822109593$ & $0.4742323085$ & $0.9811311007$\\
$1$ & $0.125$ & $25$ & $0.9950194354$ & $0.4976484589$ & $0.9950525697$\\
$6$ & $0.5$ & $38$ & $0.9930042272$ & $0.4928088283$ & $0.9928846141$\\
$6$ & $0.25$ & $75$ & $1.0019534997$ & $0.5009693722$ & $1.0019809567$\\
$6$ & $0.125$ & $151$ & $0.9982247493$ & $0.4985013113$ & $0.9982182871$\\
\bottomrule
\end{tabular}
\caption{Optimal orders use the nearest integer to $\pi x$. The bounded
rounding offset explains nonmonotonicity between rows. The last column is a
prediction, not a fitted constant.}
\label{tab:optimal}
\end{table}

The relatively small value $x=2$ in the first row is not in a very deep
asymptotic regime; the refined prediction need not improve the leading one
at every such point. At larger $x$ the explicit correction consistently
tracks the bounded-offset behavior. No monotonic rate of convergence in $h$
is claimed.

\subsection{The inverse cancellation window}
For each row of \cref{tab:window}, the script solves the scalar balance
coordinate for $x$, chooses the nearest even $M$ to $\pi x$, forms the exact
thermal-series target $y=L_h(x)$, and takes one Newton displacement for the
truncated inverse equation. It then reevaluates the \emph{nonlinear} equation
at the displaced point. The normalized residual magnitudes in the recorded
run are below $1.1\cdot10^{-44}$, $1.3\cdot10^{-70}$, and
$1.1\cdot10^{-51}$ for $h=0.5,0.25,0.125$, respectively. These are
floating-point residuals, not formal root certificates.

\begin{table}[htbp]
\centering
\begin{tabular}{@{}rrrrr@{}}
\toprule
$h$ & $s$ & $M$ & $S'(hx)\,\delta/r_h$ & $1-\ee^{-s}$\\
\midrule
$0.5$ & $0$ & $38$ & $-0.0002188003$ & $0$\\
$0.25$ & $-1$ & $78$ & $-1.7244251297$ & $-1.7182818285$\\
$0.25$ & $0$ & $78$ & $0.0031005179$ & $0$\\
$0.25$ & $1$ & $78$ & $0.6329037362$ & $0.6321205588$\\
$0.125$ & $0$ & $156$ & $-0.0006764509$ & $0$\\
\bottomrule
\end{tabular}
\caption{Normalized Newton inverse displacements in the parity-even window.
The first relative correction predicts the nonzero small values at $s=0$.}
\label{tab:window}
\end{table}

The computed centers are
\[
 x_c(0.5)=12.1852197127\ldots,\quad
 x_c(0.25)=24.6953774093\ldots,\quad
 x_c(0.125)=49.7723472861\ldots.
\]
For $h=0.25$, the refined normalized prediction at $s=0$ is
$0.0031306279$, compared with the observed $0.0031005179$.
This tests the polynomial $C_1(d)$ independently of the leading zero.

\subsection{Rational enclosures and multinomial checks}
At $t=1/2$ and $s=1/4$, a direct evaluation gives
$\D_{\log2}(1/4)=0.1354393270783473107\ldots$.
\begin{table}[htbp]
\centering
\begin{tabular}{@{}rrrr@{}}
\toprule
$N$ & $[N-1/N](1/4)$ & $[N/N](1/4)$ & Width\\
\midrule
$1$ & $0.135086455331412$ & $0.135507348193646$ & $4.20893\cdot10^{-4}$\\
$2$ & $0.135420747012517$ & $0.135445883566755$ & $2.51366\cdot10^{-5}$\\
$3$ & $0.135436670206401$ & $0.135440575785848$ & $3.90558\cdot10^{-6}$\\
$4$ & $0.135438702706998$ & $0.135439673658955$ & $9.70952\cdot10^{-7}$\\
$5$ & $0.135439129751585$ & $0.135439449154842$ & $3.19403\cdot10^{-7}$\\
\bottomrule
\end{tabular}
\caption{Decimal displays of exact rational Pad\'e bounds. The JSON output
contains the full rational numbers; rounded table entries themselves are not
outward-rounded enclosures.}
\label{tab:pade}
\end{table}

For $h=0.2$ and $x=10$, the ratio of the positive multinomial remainder to
the shortest-block prediction is $1.0053351399$ for both $(1,2,3)$ and
$(1,1,2)$, with orders $M=31$, and $0.9958300329$ for $(2,2,3)$, with
$M=63$. The first two ratios agree to the displayed precision although their
shortest-block multiplicities differ; the predicted multiplicity was divided
out before comparison.

\subsection{How to choose an order without knowing the inverse}
In the compact crossover regime, first solve the zeroth-order composite
 equation $B_0(h,z_0)=y$ locally. The exact residual is
$\E(h)+R_0=O_K(h)$, while its derivative is bounded below, so
$z_0-x=O_K(h)$. Choosing $M$ to be the nearest integer of the desired parity
to $\pi z_0$ therefore ensures $|M-\pi x|=O(1)$, as required by the sharp
 theorem. This is stronger than using only the phase center, whose missing
logarithmic amplitude would cause an $O(|\log h|)$ index displacement.

For ordinary accuracy, solve with the modular completion retained and use
signed or Pad\'e enclosures. Deliberately omitting it to exploit cancellation
is a specialized diagnostic, not the default recommendation for a robust
numeric inverse. A certified implementation should adapt the order using the
actual remainder enclosure rather than rely on asymptotic balance alone.

\section{Repository integration and proof-assistant boundaries}
\label{sec:integration}

The inspected snapshot is the full commit \texttt{\commit}. The source
chapter is in
\begin{quote}\small\ttfamily
Analysis/FabiusFunction/docs/\\
semi-formalized-research-frontiers/drafts/\\
series-and-transseries/\\
Combinatorial\_Transseries\_Inverses/\\
Combinatorial\_Transseries\_Inverses.tex.
\end{quote}
% ed. (2026-09-29): current path added by ProveIt.
\noindent Editorial note (ProveIt, 2026-09-29): this is the pre-split path;
the companion is now
\begin{quote}\small\ttfamily
Analysis/Transseries/docs/series-and-transseries/\\
Combinatorial\_Transseries\_Inverses/\\
Combinatorial\_Transseries\_Inverses.tex.
\end{quote}
The following crosswalk prevents a formal coefficient identity from being
silently upgraded into an analytic theorem.

\begingroup\small
\begin{longtable}{@{}>{\raggedright\arraybackslash}p{.24\textwidth}>{\raggedright\arraybackslash}p{.34\textwidth}>{\raggedright\arraybackslash}p{.35\textwidth}@{}}
\toprule
Repository item & Already present there & Added here\\
\midrule
\endhead
\texttt{t3:prop:double-A}
& Fixed-order tail expansion for $\tau\ge\tau_0>0$.
& Exact sign, explicit next-term majorant, and uniform endpoint estimate
\eqref{eq:uniform}.\\[4pt]
\texttt{t3:eq:modularP} and \texttt{t3:eq:modularsectors}
& Exact modular identity and divisor-sum completion.
& Quantitative competition with the optimal finite-size error; no claim to a
new modular identity.\\[4pt]
\texttt{t3:eq:Sphase} and \texttt{t3:eq:Sderivative}
& Dilogarithmic phase and elementary positive derivative.
& Uniform lower bound for the \emph{exact} inverse slope and exponentially
accurate residual transport.\\[4pt]
\texttt{t3:eq:crossover-all}
& All formal finite-order inverse coefficients and fixed-order asymptotics.
& Sharp optimal-order remainder, parity profile, and exact cancellation
location.\\[4pt]
Canonical inverse certification framework
& Residual brackets and order-theoretic threshold principles, with stated
formalization scopes.
& Explicit signed and rational analytic forward envelopes suitable as inputs
to that framework.\\
\bottomrule
\end{longtable}
\endgroup

A useful formalization can be modular. The finite identity in
\cref{lem:thermal} after the coth expansion is supplied is algebraic.
Positivity and finite-cutoff monotonicity of the remainder kernel are
straightforward ordered-field statements. Summability and passage to the
infinite sums form the analytic layer. The residual transport theorem then
uses only a derivative lower bound and a root-existence bracket.

The hardest new formal components are the bilateral large-moment identity,
its uniform exponential estimate, the moment expansion with a bounded integer
offset, and the quadrature construction. They should not be bundled into one
opaque ``transseries theorem''. The repository's existing distinction between
formal identities, conditional analytic lemmas, and assembled asymptotic
statements is the appropriate design discipline \cite{repo-readme}.

No repository files were modified for this deliverable. The present proofs
were not compiled in Lean, and no external repository theorem was accepted
solely because a README called it proved. The argument uses the printed
classical identities and reproduces the mathematical derivations needed for
its own claims.

\section{Further research questions and concrete next targets}
\label{sec:future}

The following are proposed problems, not results claimed in the present
paper. Each is narrow enough to expose a verifiable output.

\begin{enumerate}[label=\textbf{Q\arabic*.},leftmargin=13mm]
\item \textbf{Sharp optimal remainders uniformly down to $hx=0$.}
The fixed-order bound is already uniform there. Can the prefactor and
$C_1(d)$ in \cref{thm:optimal} be made uniform for $0<hx\le\tau_1$, assuming
only $x\to\infty$? The exact power-sum identity suggests that a strengthened
version should be possible, but all differentiated and discarded-weight
estimates must be uniform simultaneously.

\item \textbf{The discrete-saddle transition at large $hx$.}
The relevant moment distribution has center approximately $2\pi/h$ and
width approximately $\sqrt{2\pi x}/(hx)$. Analyze the transition where this
width is bounded, rather than large. A theta-modulated prefactor, depending
on the fractional position of the saddle, is a concrete candidate. This
would connect the present double-scaling theorem to fixed-$q$ asymptotics
without an unjustified continuum replacement.

\item \textbf{Higher cancellation corrections and parity optimization.}
Compute the $x^{-2}$ and $x^{-3}$ relative corrections to \eqref{eq:optimal},
with uniform remainders, and propagate them to $x_{h,M}$. The dependence on
integer rounding must remain explicit. A practical goal is a certified rule
selecting the best of two nearby even orders from the first few correction
polynomials.

\item \textbf{Coalescing shortest blocks.}
For multinomials with $\alpha_2-\alpha_1=O(1/x)$, a second block is no longer
exponentially negligible in the large moment. Derive a uniform replacement
for the discrete multiplicity $m_*$, with a continuous transition profile
in $x(\alpha_2-\alpha_1)$. Fixed block proportions in \cref{thm:shortest}
do not cover this problem.

\item \textbf{Which signed $q$-product quotients have a positive kernel?}
For a general finite quotient, the weight is an exponential polynomial
$\sum_jc_j\ee^{-a_j u}$. Characterize useful, algorithmically testable
conditions ensuring nonnegativity for every $u>0$. Majorization and Laplace
transform ordering are possible approaches, but neither is assumed to give
a complete criterion.

\item \textbf{Quantitative Pad\'e convergence in joint limits.}
\Cref{thm:pade} proves convergence for fixed $\tau,s$. Obtain computable
rates that remain useful as $\tau\to0$, $s\to0$, and $N\to\infty$ together.
The measure has an unbounded support and factorial moments, so compact-support
Pad\'e estimates cannot be imported unchanged.

\item \textbf{An interval-arithmetic inverse and exact integer thresholds.}
Implement outward-rounded evaluation of the signed kernel or Pad\'e bounds,
including modular-product tails. Prove termination of a bracket-refinement
algorithm away from exact thresholds, and use exact finite products to
resolve the equality cases. The supplied script is a test harness, not this
certifier.

\item \textbf{Complex $h$, actual Stokes data, and inverse resurgence.}
Determine the uncancelled divisor of \eqref{eq:poleset}, compute lateral
continuations where necessary, and propagate their differences through
compositionally inverted functions. Such a theory should explicitly
separate Stokes jumps from the real truncation-error cancellation established
here.

\item \textbf{Normalization and flat-completion uniqueness.}
The algebraic series cannot by itself determine the modular term.
Identify a minimal combination of a $q$-difference equation, growth
conditions, and a boundary normalization that uniquely recovers the chosen
Gaussian interpolation among functions sharing the same algebraic expansion.

\item \textbf{Shifted quotients such as $q$-Catalan families.}
Block lengths of the form $\alpha_i x+\beta_i$ introduce fixed shifts in the
exponential weights. Determine the leading optimal remainder and cancellation
window including the shift-dependent prefactor. This is not obtained by
simply replacing $x$ with $x+1$ in the central-binomial answer.

\item \textbf{Theta-weighted finite-field enumeration.}
Extend the positive-kernel viewpoint to the theta-modulated Galois and flag
families already studied in the companion. Establish precisely when a
positive measure survives the additional lattice sum, and when oscillatory
or signed kernels require a different certification method.

\item \textbf{A composable Lean certificate library.}
Formalize a finite geometric remainder interface, positive summation bounds,
residual-to-root transport, and a symbolic truncation checker before attempting
full Borel or quadrature formalization. A concrete first milestone is a
machine-checked finite-cutoff version of \eqref{eq:exactR} with an explicit
analytic tail supplied as a hypothesis.
\end{enumerate}

\section{Conclusion}

The singular $q\to1$ transition is not adequately described by a formal
inverse expansion plus an isolated modular term. Its finite-size Bernoulli
series has a structured, signed remainder of its own. A positive thermal
kernel makes that structure explicit and gives three kinds of control:
fixed-order bounds uniform toward the classical endpoint, rational
Stieltjes enclosures, and sharp optimal-order asymptotics.

The central quantitative conclusion is the stable inverse-error law
\[
 \widehat x_M-x=
 \frac{\ee^{-4\pi^2/h}+(-1)^{M+1}\ee^{-2\pi x}/(\pi\sqrt{x})}
      {2\log(1+\ee^{hx})}
 +O\!\left(h\left\{\ee^{-4\pi^2/h}
              +\frac{\ee^{-2\pi x}}{\pi\sqrt{x}}\right\}\right),
\]
under the precisely stated compact double-scaling hypotheses. Its balance
coordinate includes a logarithmic prefactor shift, its sign depends on
truncation parity, and its refined cancellation location is explicit.
The multinomial version shows that the shortest block and its multiplicity
control the leading finite-size error, while the full composition controls
the inverse slope.

These are delimited analytic refinements of an existing transseries program.
They do not settle general questions about transseries fields or every
$q$-product regime. They do turn an explicitly unresolved error-analysis
boundary in the inspected chapter into a collection of proved formulas and
testable, reusable inversion tools.

\appendix
\section{Explicit bounds for finite computation}
\label{app:tails}

These estimates make the positive formulas numerically usable without
pretending that a finite sum is an infinite one. Let $p=2M$, let the positive
weight satisfy $0<w_m\le b\ee^{-\lambda m}$, and choose an integer
$N\ge p/\lambda$. For a multinomial, $\lambda=\alpha hx$; in the central
case $b=2$ and $\lambda=hx$.

Since $u^p\ee^{-\lambda u}$ decreases for $u\ge N$,
\begin{equation}
 \sum_{m>N}m^p\ee^{-\lambda m}
 \le\int_N^\infty u^p\ee^{-\lambda u}\dd u
 =\frac{\Gamma(p+1,\lambda N)}{\lambda^{p+1}}.
 \label{eq:mtail}
\end{equation}
Here $\Gamma(a,z)=\int_z^\infty u^{a-1}\ee^{-u}\dd u$ is the upper incomplete
gamma function, a finite exponential-polynomial expression when $a=p+1$
is a positive integer. The corresponding absolute remainder-tail bound is
\begin{equation}
 E_m\le
 \frac{h\zeta(p+2)}{2\pi^2}
      \left(\frac h{2\pi}\right)^p
      \frac{b\Gamma(p+1,\lambda N)}{\lambda^{p+1}}.
 \label{eq:Rmtail}
\end{equation}
For the omitted $k>K$ terms at retained $m\le N$, write
$z_m=(hm/(2\pi))^2$ and $v_m=w_m z_m^M$. Then
\begin{equation}
 E_k\le\frac h{2\pi^2}
           \frac{\sum_{m\le N}v_m}{(p+1)K^{p+1}}.
 \label{eq:Rktail}
\end{equation}
This is the integral test for $\sum_{k>K}k^{-p-2}$. The positive rectangular
partial sum of \eqref{eq:exactR}, without its parity sign, is a lower bound;
adding $E_m+E_k$ is an upper bound. No cancellation of huge numbers is needed
for the small remainder itself.

For the independent thermal reference \eqref{eq:thermalreference}, the
bracketed kernel lies between zero and $1/2$, since
$\ee^z-1>z$ and its cotangent representation is positive after the first two
terms are removed. Therefore the $m>N$ tail is at most
\begin{equation}
 \frac{b\ee^{-\lambda(N+1)}}{2(N+1)(1-\ee^{-\lambda})}.
 \label{eq:Ftail}
\end{equation}
Finally, if the modular divisor sum is stopped after $J$ terms,
\begin{equation}
 0<\E(h)-\sum_{n=1}^{J}\sigma_{-1}(n)r_h^n
 \le\frac{r_h^{J+1}\{(J+1)-Jr_h\}}{(1-r_h)^2}.
 \label{eq:Etail}
\end{equation}
All these are analytic inequalities. Outward rounding is still necessary
when implementing their right-hand sides as rigorous numerical enclosures.

\section{Coefficient construction and checks of normalization}
\label{app:coefficients}

The moments needed by the Pad\'e construction can be evaluated without
symbolic differentiation of large rational expressions. For $n\ge0$,
\begin{equation}
 \Li_{-n}(t)=\sum_{j=0}^{n}j!\left\{\begin{matrix}n+1\\j+1\end{matrix}\right\}
                         \left(\frac t{1-t}\right)^{j+1},
 \label{eq:stirlingpolylog}
\end{equation}
where braces denote Stirling numbers of the second kind. To prove the
identity, start with $\Li_0(t)=t/(1-t)$, apply $t\,d/dt$, and use the Stirling
recurrence. At $t=1/2$, \eqref{eq:moments} becomes the explicitly rational
finite sum
\begin{equation}
 \mu_j=\frac{|B_{2j+2}|}{(2j+2)!}
 \sum_{k=0}^{2j}k!\left\{\begin{matrix}2j+1\\k+1\end{matrix}\right\}
                    \left(2-3^{-(k+1)}\right).
 \label{eq:rationalmoments}
\end{equation}
This is the formula executed in the script.

For a formal series $f(s)=\sum_{j\ge0}c_js^j$, a normalized Pad\'e denominator
$q(s)=1+\sum_{j=1}^Nq_js^j$ for numerator degree $L$ is found from
\begin{equation}
 \sum_{j=1}^{N}q_jc_{k-j}=-c_k,\qquad k=L+1,\ldots,L+N,
 \label{eq:padelinear}
\end{equation}
with $c_j=0$ for $j<0$. The numerator coefficients are
$p_k=\sum_{j=0}^{\min(k,N)}q_jc_{k-j}$. In our application
$c_j=(-1)^j\mu_j$. Positive definiteness of the moment form gives the
nondegeneracy needed in \cref{thm:pade}; the program additionally checks the
relevant finite determinants exactly.

Several normalization traps are worth recording. The Bernoulli sum in
$B_M$ contains $M$ terms, whereas the first omitted moment has power $2M$.
The remainder sign is $(-1)^{M+1}$, not $(-1)^M$, because the tail enters as
$A_h(t^2)-2A_h(t)$. The next-term leading prefactor is \emph{twice} the
actual optimal remainder prefactor. The modular contribution to $L_h$ is
positive, because it is $-\log\qpoch{r_h}{r_h}$. Finally, the inverse
approximation omitting that contribution is shifted by $+\E/S'$ to leading
order; interchanging the exact and approximate inverse reverses the sign.

\section{Package contents and build instructions}
\label{app:build}

The archive is self-contained and does not require the repository's shared
notation file. It contains the editable \texttt{q\_transition.tex}, its
compiled PDF, this manuscript's verification script and recorded JSON output,
a minimal dependency file, a provenance record, and a README.

Compile the document with a current TeX Live installation containing
\texttt{newpx}, \texttt{tcolorbox}, \texttt{hyperref}, and \texttt{cleveref}:
\begin{verbatim}
pdflatex -interaction=nonstopmode -halt-on-error q_transition.tex
pdflatex -interaction=nonstopmode -halt-on-error q_transition.tex
pdflatex -interaction=nonstopmode -halt-on-error q_transition.tex
\end{verbatim}
The bibliography is embedded, so BibTeX is unnecessary. To reproduce the
checks, install the versions in \texttt{verification/requirements.txt} and run
\begin{verbatim}
python verification/verify.py
\end{verbatim}
The program overwrites only its local \texttt{verification/results.json}.
It makes no network request and does not alter the user's repository.

\begin{thebibliography}{99}
\raggedright
\bibitem{repo}
V. Reshetnikov, \emph{ProveIt}, repository snapshot
\texttt{\commit}. Companion manuscript,
\emph{Combinatorial Transseries and Their Inverses}, chapter
``The singular transition $q\to1$'', labels
\texttt{t3:prop:double-A}, \texttt{t3:eq:modularP},
\texttt{t3:eq:Sphase}, \texttt{t3:eq:crossover-all}.
Inspected 29 September 2026.
\url{https://github.com/VladimirReshetnikov/ProveIt/tree/20301601562cf84173f2b503b0dc9c3fbf21cd0c}.

\bibitem{repo-readme}
V. Reshetnikov, \emph{ProveIt}, ``Series and transseries'' and
``Transseries: the polynomial--logarithmic calculus and its inversions'',
repository README files at the same snapshot. The source inventory explicitly
separates formal identities, partial analytic crosswalks, and absent assembled
formalizations. See the source path in \cref{sec:integration}.

\bibitem{dlmf-hyper}
NIST Digital Library of Mathematical Functions,
\S4.36, ``Infinite Products and Partial Fractions'', especially equations
4.36.3 and 4.36.4. Accessed 29 September 2026.
\url{https://dlmf.nist.gov/4.36}.

\bibitem{dlmf-eta}
NIST Digital Library of Mathematical Functions,
\S23.18, ``Modular Transformations'', especially equation 23.18.5.
Accessed 29 September 2026.
\url{https://dlmf.nist.gov/23.18}.

\bibitem{zhang}
C. Zhang, ``A modular-type formula for $(x;q)_\infty$'',
\emph{The Ramanujan Journal} \textbf{46} (2018), 269--305.
\url{https://doi.org/10.1007/s11139-017-9967-5}.

\bibitem{fantini-rella}
V. Fantini and C. Rella, ``Modular resurgence, $q$-Pochhammer symbols, and
quantum operators from mirror curves'', \emph{Letters in Mathematical Physics}
\textbf{116} (2026), article 39.
\url{https://doi.org/10.1007/s11005-026-02063-x}.
Preprint \url{https://arxiv.org/abs/2506.08265}, version 2, 1 April 2026.

\bibitem{simon}
B. Simon, ``The classical moment problem as a self-adjoint finite difference
operator'', \emph{Advances in Mathematics} \textbf{137} (1998), 82--203.
\url{https://arxiv.org/abs/math-ph/9906008}.
% ed. (2026-09-29): editorial bibliography entries added by ProveIt.
\bibitem{ed:siblings}
[Editorial addition, ProveIt, 2026-09-29.] The five packages of the
$q\to1$ merge unit in the ProveIt series-and-transseries collection,
\path{Analysis/Transseries/docs/series-and-transseries/}:
\path{Uniform_q_Multinomial_Certified_Inversion/uniform_q_multinomial_transseries.tex};
\path{Uniform_Resurgent_Crossover_Gaussian_Binomials/article.tex};
\path{Certified_Inversion_q_to_1_Transition/q_transition.tex};
\path{Gamma_Core_q_to_1_Crossover/Gamma_Core_Transseries.tex};
\path{Theta_Resolved_Optimal_Truncation_q_Multinomial/article.tex}.
Filed in batches 45 and 46 of \path{docs/incoming/README.md}; unrefereed.

\bibitem{ed:gcc}
[Editorial addition, ProveIt, 2026-09-29.] ProveIt, \emph{Gaussian
Coefficient Calculus},
\path{Analysis/FabiusFunction/docs/semi-formalized-research-frontiers/drafts/combinatorial-coefficient-calculus/Gaussian_Coefficient_Calculus/Gaussian_Coefficient_Calculus.tex},
Theorem \texttt{q3:thm:double-scaling} (line 3942) and its domain
\texttt{q3:eq:double-scaling-domain}.

\bibitem{ed:tai}
[Editorial addition, ProveIt, 2026-09-29.] ProveIt, \emph{Transseries: the
polynomial-logarithmic calculus, series reversal at infinity, and
inversion} (canonical volume),
\path{Analysis/Transseries/docs/series-and-transseries/Transseries_And_Inversion/transseries_and_inversion.tex},
Theorem \texttt{p0:thm:optimal-truncation} (line 39360).

\bibitem{ed:lean}
[Editorial addition, ProveIt, 2026-09-29.] ProveIt Lean library:
\path{Fabius.exists_eq_in_residual_interval} in
\path{Analysis/FabiusFunction/Lean/FabiusFunction/MeanValueBracket.lean}
and \path{Fabius.transport_bound} (the crosswalk of the canonical volume's
\texttt{p0:eq:transport-bound}) in
\path{Analysis/FabiusFunction/Lean/FabiusFunction/RemainderTransport.lean}.
\end{thebibliography}
\end{document}
